\subsection{Passage on the no-root side of the fold}
\label{sec:direct-no-root-fold}

Suppose now that the exact minimum of \(K_+\) in
\eqref{eq:direct-fold-region} is positive. The normalized drift
\(K_+\) is then positive on the whole positive unscaled radius
interval: outside the fixed fold collar its quadratic
part has a positive margin, and its cubic term is positive.
As this minimum tends to zero, we need estimates uniform in the
possibly divergent bottleneck action. The bounded-output result
compares half-orbits whose projected states at a fixed secondary
height stay bounded. For larger projected outputs, we isolate the
bottleneck action in the exact exponential transport, prepare the
outgoing height uniformly, and compare exits on the primary and
cubic scales.

\begin{proposition}[Zero bound across the fold]
\label{prop:direct-negative-fold}
For fixed sufficiently small \(\epsilon_{\rm f}>0\), \(\Delta_p\) has at most
three distinct zeros throughout \eqref{eq:direct-fold-region}, including
the fold and both its sides, for all sufficiently small independent
\(s,r>0\).
\end{proposition}

\subsubsection{Bottleneck transport and escaped central heights}

Fix the postfold cut \(K_0>2\) from the annular construction,
and later a larger fixed \(M>K_0\). A central height is called
escaped if its positive half-orbit is still thin at \(x=M\):
\begin{equation}
\label{eq:direct-fold-escaped}
 z(Msr;h)/(s^3M^2r^2)\le\delta_{\rm an} .
\end{equation}
This is one interval of central heights.
Use a single boundary-flow coordinate on a fixed positive
box from \(x_L=x_{\rm f}-2\epsilon_{\rm f}\) to beyond \(K_0\).
The bottleneck formula \eqref{eq:direct-bottleneck-action}
gives
\[
 \Psi_{K_0}(y_{K_0})
    =\e^{\mathcal A_f/s}\Psi_L(y_L),\qquad
 \mathcal A_f=\int_{x_L}^{K_0}\frac{x\,dx}{K_+(x)} .
\]
Here \(Q=s^3r^2\) and \(y_j=z(jsr;h)/Q\) for
\(j=x_L,K_0\).
The endpoint coordinates have uniform bounds because their integrals
of \(1/F_+\) lie on regular collars. The full bottleneck action remains
in the explicit exponential factor; the composition below is
differentiated in the expanded height, where the inverse bounds stay
uniform.
The construction following \eqref{eq:direct-bottleneck-action} applies
with \(x_R=K_0\), \(\delta=s\), \(\beta=\beta_f\), and the regular
central prefix of \cref{lem:direct-source-height}, of action
\(\mathcal A_L=\int_0^{x_L}x/K_+(x)\,dx\), where
\[
 \mathcal A_{K_0}=\mathcal A_L+\mathcal A_f,\qquad
 \beta_f=\e^{-\mathcal A_f/s},\qquad
 \xi=\e^{\mathcal A_{K_0}/s}h/Q .
\]
By \eqref{eq:direct-bottleneck-inverse-composition},
\begin{equation}
\label{eq:direct-fold-inverse-composition}
 y_{K_0}=\xi V_p(\xi),\qquad V_p(0)=1,
\end{equation}
on a fixed small \(\xi\)-interval whose image contains a fixed small
output interval; the unit, its reciprocal, and the map inverse have
uniformly bounded \(\xi\)-derivatives.
Composing with the annular coordinate gives
\eqref{eq:direct-annulus-prefix} with action \(\mathcal A_{K_0}\).

We next check that every escaped central height lies in the fixed
domain of this inverse construction. For such a height, put
\(y_a(\rho)=z(s\rho;h)/(s^3\rho^2)\).
For fixed \(M>K_0\) and sufficiently small \(r\),
the segment \(K_0r\le\rho\le Mr\) lies in the
positive annular box, where \eqref{eq:direct-annulus-level-crossing}
excludes a downward crossing of \(y_a=\delta_{\rm an}\).
The escaped condition at \(Mr\) implies
\(y_a\le\delta_{\rm an}\) on the whole preceding segment.
The annular transport \eqref{eq:direct-annulus-action} yields
\[
 \begin{gathered}
 \frac{\widehat z(K_0r)}{Q}
   =\e^{-\mathcal J_{i,M}/s}\frac{\widehat z(Mr)}{Q}
   \le CM^2\delta_{\rm an}\e^{-\mathcal J_{i,M}/s},\\
 \mathcal J_{i,M}=\int_{K_0r}^{Mr}
       \frac{dt}{tK_+(r/t,t)}
       \ge c\log(M/K_0)>0.
 \end{gathered}
\]
At the inner cut, \(y_{K_0}=K_0^2y_a(K_0r)\) is
uniformly comparable to \(\widehat z(K_0r)/Q\).
Thus \(y_{K_0}\) lies in the small output domain of the bottleneck
and central-prefix inverse coordinates. Their composition gives a
small \(\xi\), and flow uniqueness identifies the recovered central
height with \(h\). Hence
\eqref{eq:direct-annulus-prefix} applies throughout the interval of
escaped central heights.

For every escaped central height, the positive half-orbit crosses
the fold collar below a bounded secondary height. The action gap gives
\begin{equation}
\label{eq:direct-escaped-source-gap}
 h/(s^3r^2)\le C_{M,\delta_{\rm an}}\e^{-\mathcal A_{\rm gap}/s},\qquad
 \mathcal A_{\rm gap}\ge c/\epsilon_{\rm f}.
\end{equation}
Fix the negative secondary cuts needed below, then choose
\(\epsilon_{\rm f}\) small so that \(\mathcal A_{\rm gap}\) exceeds their
regular actions by a fixed gap.
At each such cut the expanded negative input is exponentially small
by \eqref{eq:direct-escaped-source-gap}. The regular inverse
central-height coordinate therefore constructs a thin negative prefix;
flow uniqueness identifies it with the branch from the original \(h\).
Both signs now satisfy the hypotheses of
\cref{lem:direct-annulus-rates} whenever their exits
have small primary radius, and the negative branch satisfies those of
\cref{lem:negative-branch-far-root}.

\subsubsection{Both exits on the small primary annulus}

Take small fixed annular margins \(u_0,v_0>0\),
with \(u_0\) eventually much smaller than \(v_0\).
Choose \(M\) and the action gap so that both inner ratios
\(u_\sigma=r/\rho_\sigma\) are at most \(u_0\).
If both thin exits satisfy \(v_\sigma=\rho_\sigma\le v_0\), the prefixes
constructed above meet \cref{cor:direct-annular-pair}. Since \(b\ge1/2\),
small annular margins followed by small \(s\) give \(\mathcal L''<0\)
in \(T=-s\log h\).

\subsubsection{The negative primary and far-root outputs}

When the positive exit is at most \(v_0\), its annular
rate satisfies
\[
 \mathfrak H_+\ge
       2-Cu_0-Cv_0^2-Cs(1+|\log s|).
\]
For a negative exit in \(v_0/2\le\rho_-\le0.9\),
\eqref{eq:direct-negative-primary-rates} gives
\(\mathfrak H_-=2-3\rho_-+o(1)\).
First choose \(v_0\) small, then \(u_0\) much smaller
than \(v_0\), and finally reduce \(r,s\).
The rate difference \(\mathfrak H_--\mathfrak H_+\) is strictly negative.
The same sign
holds if the negative graph is in its far-root interval,
by \eqref{eq:direct-negative-far-rate}.
The exterior contribution to \(\mathcal L'\) tends to
zero with the negative endpoint derivative.
Hence \(\mathcal L'<0\) on both ranges.

It remains to locate the negative branch when the positive exit has
left the small annulus. For an escaped central height, a positive
annular thin exit with \(\rho_+\ge v_0\) places the negative graph in
the far-root interval, as the following comparison of central heights
shows.
Choose first the compact negative comparison cut \(0.9\) and the
annular cut \(v_0\), and let \(C_{\rm th}\) bound all fixed endpoint and
prefix terms in the two central-height equations.  Then reduce the fold
width \(\epsilon_{\rm f}\), after the preceding choices, so that
\(\mathcal A_{\rm gap}>2C_{\rm th}\).  Compare the central height for
which the positive exit is at \(v_0\) with the central height for which
the negative height at \(0.9\) equals the thin threshold.  Their actions satisfy
\begin{equation}
\label{eq:direct-fold-far-threshold}
 \begin{split}
 \mathcal A_+(v_0)
    &=\mathcal A_{K_0}+\log\frac{v_0}{K_0r}
                                   +O(K_0^{-1}+v_0),\\
 \mathcal A_-(0.9)&=\log(1/r)+O(1).
 \end{split}
\end{equation}
The first estimate follows by integrating the annular factor
\(K_+(r/\rho,\rho)^{-1}-1=O(r/\rho+\rho)\) after \(K_0r\).
For the second, the negative secondary prefix has no
bottleneck, and the remaining primary action has fixed
endpoint margins.
Since \(\mathcal A_{K_0}\ge\mathcal A_{\rm gap}\), the displayed estimates
and the order of choices give, uniformly for small \(r\),
\begin{equation}
\label{eq:direct-fold-far-action-gap}
 \mathcal A_+(v_0)-\mathcal A_-(0.9)
 \ge \mathcal A_{\rm gap}-C_{\rm th}\ge C_{\rm th}>0.
\end{equation}
Writing \(h_+^{v_0}\) and \(h_-^{0.9}\) for the two central heights
(with \(h_-^{0.9}\) defined by the thin threshold at the fixed radius
\(0.9\)),
\eqref{eq:direct-annulus-source} and the uniform boundedness of its
height factors give
\begin{equation}
\label{eq:direct-fold-far-height-order}
 s\log\frac{h_+^{v_0}}{h_-^{0.9}}
 =-\bigl[\mathcal A_+(v_0)-\mathcal A_-(0.9)\bigr]+O(s)<0
\end{equation}
after reducing \(s\).  Thus \(h_+^{v_0}<h_-^{0.9}\).  The exit-height
maps are strictly decreasing in their radii by the event equation, so
for any \(\rho_+\ge v_0\) its central height satisfies
\(h\le h_+^{v_0}<h_-^{0.9}\). The negative graph at that same \(h\)
is consequently still thin at \(0.9\): by
\cref{lem:central-height-order}, its height at this fixed radius is
smaller than the thin threshold attained by \(h_-^{0.9}\).
Part (ii) of \cref{lem:negative-branch-far-root} now gives
\(w_n>0.89\), hence \eqref{eq:direct-negative-far-core}.

\subsubsection{Positive continuation through the cubic regime}

For escaped central heights whose positive annular exit is at least
\(v_0\), which form an interval, we apply
\cref{lem:positive-branch-adaptive-exit}. Its prefix hypothesis
\eqref{eq:positive-adaptive-prefix} is
\eqref{eq:direct-fold-inverse-composition}, with \(\mathcal V=V_p\) and
\(\mathcal A_i=\mathcal A_{K_0}\). Escaped central heights are thin
at \(K_0r\), and their inverse estimate makes \(\xi\) uniformly small,
so it lies in the domain of \(V_p\) after reducing \(s\). By
\eqref{eq:positive-annular-adaptive-heights}, the adaptive exit follows
the annular thin exit; hence \(\varrho_e\ge sv_0\). By the first estimate in
\eqref{eq:direct-fold-far-threshold},
\(\mathcal J_+(v_0)=\mathcal A_+(v_0)-\mathcal A_{K_0}\ge\log(1/r)-C\), which is at
least \(\frac12\log(1/r)\) after fixing a sufficiently small upper bound
for \(r\). Thus part (iv) of the lemma applies with \(\rho_{\rm g}=v_0\)
and \(c_{\rm g}=1/2\). It gives \eqref{eq:direct-no-root-adaptive-rate},
\(\mathfrak H_+>c\varrho_e^+/s\), at these exits. The bounds on \(V_p\)
in \eqref{eq:direct-bottleneck-inverse-composition} make this estimate
uniform in the bottleneck action.
The negative graph is in the far-root interval by
\eqref{eq:direct-fold-far-threshold}. Its strictly
negative rate and the small exterior term imply
\(\mathcal L'<0\) here as well.

We now combine the exit comparisons on the interval of escaped heights.
On the interval where both exits are at most \(v_0\), we have
\(\mathcal L''<0\).
Within the escaped central heights, the condition
\(y_-(v_0/2;h)\le\delta_{\rm an}\) defines a second interval,
because the height at a fixed radius increases with \(h\).
The negative primary, negative far-root, and positive adaptive
estimates cover this interval, including their fixed overlaps.
Thus \(\mathcal L'<0\) throughout the second interval.
These two intervals cover
\eqref{eq:direct-fold-escaped}: a positive exit at least
\(v_0\) forces the negative far-root interval; otherwise
either both exits are small or the negative graph is
already in the second interval.
The estimates for \(\mathcal L(T)\) include both cuts
\(v_0/2\) and \(v_0\).
Consequently, on the entire escaped interval,
\begin{equation}
\label{eq:direct-escaped-critical-maximum}
 \mathcal L'=0\quad\Longrightarrow\quad\mathcal L''<0.
\end{equation}

\subsubsection{Covering the height interval}

Set \(B_*=4M\) in \cref{prop:direct-fold-bounded}.
Its bounded-output estimates give
\eqref{eq:direct-fold-critical-maximum}.
The quadratic-prefix lemma applies on the early interval
\[
 h\ge z_e=s^3r^2\delta_e,\qquad\hbox{or}\qquad
 h<z_e,\quad x_+(z_e;h)\le x_{\rm f}-\epsilon_{\rm f}/2 .
\]
Its regular-prefix hypotheses hold on the no-root side. In \(T\),
the early heights form an initial interval, and each zero of \(\mathcal L\) on it
satisfies \(\mathcal L'>0\).

For every other central height, the projected state of the positive
half-orbit at its hit of \(x=x_{\rm f}-\epsilon_{\rm f}/2\) is greater than
\(x_{\rm f}-3\epsilon_{\rm f}/4\).
Follow the half-orbit to the first event \(Z=Z_f\) or \(x=3M\).
If the height is reached first, the projected output
lies between \(x_{\rm f}-\epsilon_{\rm f}\) and \(B_*\), and is counted by
the bounded-output proposition.
If \(x=3M\) is reached first, the central height satisfies
\eqref{eq:direct-fold-escaped}. Indeed, between \(x=M\) and that event
the half-orbit lies in the coordinate box of
\cref{sec:direct-negative-fold}, which extends beyond \(B_*=4M\), and
\(B_0Z_f\le M^2\) since \(Z_f\) is small. For small \(s\),
\cref{lem:direct-radial-advance} with \(\delta=s\), \(X_1=M\), \(X_2=3M\)
and \(C_F=\sup|F_+|\) on that box bounds the height at \(x=M\) by
\(Z_f\e^{-M/(C_Fs)}<sM^2\delta_{\rm an}\).
Every admissible half-orbit under consideration reaches one of these
events before its fixed endpoint.

The complement of the early set is the single interval
\[
 J=\{T:h\in\cH_p,\ h<z_e,\quad
                      x_+(z_e;h)>x_{\rm f}-\epsilon_{\rm f}/2\}.
\]
Every point of \(J\) satisfies a bounded-output or escaped estimate.
Their common conclusion is the critical-point inequality
\(\mathcal L'=0\Rightarrow\mathcal L''<0\), by
\eqref{eq:direct-fold-critical-maximum} and
\eqref{eq:direct-escaped-critical-maximum}, for the same smooth function
on \(J\). The early cut
belongs to the initial interval, where
\eqref{eq:direct-quadratic-prefix-zero-orientation} includes the
terminal radius. Thus the initial interval and \(J\) cover the
whole height interval \(\cH_p\).
Apply \cref{cor:ordered-cover-zero-count}, with \(J\) as the
middle piece and an empty final piece. It gives the bound of three on
the no-root side. The same bound on the
root-present side follows from \cref{prop:direct-fold-bounded}.
